\documentclass[11pt,a4paper]{article}
\usepackage[T1]{fontenc}
\usepackage[slovene]{babel}
\usepackage{amsthm}

\newtheorem{izrek}{Izrek}[section]

\title{Trikotniki}
\author{Matija Pretnar}
\date{}

\begin{document}
\maketitle

\section{Uvod}
Trikotnike poznamo že iz osnovne šole.

\section{Pitagorov izrek}
\begin{izrek}[Pitagorov izrek]
  \label{izr:pitagora}
  V pravokotnem trikotniku s katetama~\(a\) in~\(b\) ter hipotenuzo~\(c\) velja
  \(a^2 + b^2 = c^2\).
\end{izrek}

Iz izreka~\ref{izr:pitagora} izpeljemo razdaljo med točkama v prostoru~\cite{vidav08}.

\bibliographystyle{plain}
\bibliography{literatura}

\end{document}
